\section*{Hoofdstuk \ref{ch:b2:cell-signalling} --- Chemische boodschappers en signaaltransductie}

\begin{solution}{exo:b2:cell-signalling:1}
Insuline: endocrien, wateroplosbaar (peptide). Cortisol: endocrien,
vetoplosbaar. Acetylcholine in een synaps: synaptisch, wateroplosbaar.
Histamine in een wond: paracrien, wateroplosbaar. Oestradiol: endocrien,
vetoplosbaar. Adrenaline uit de bijnier: endocrien, wateroplosbaar.
Stikstofmonoxide: paracrien, vetoplosbaar (een gas dat membranen
passeert).
\end{solution}

\begin{solution}{exo:b2:cell-signalling:2}
Adrenaline bindt de receptor (uit: dissociatie, internalisatie van de
receptor); de receptor activeert het G-eiwit, GDP voor GTP (uit: hydrolyse van
GTP in seconden); het G-eiwit activeert adenylaatcyclase, dat cAMP maakt
(uit: fosfodiësterase); cAMP activeert kinase A (uit: verlies van cAMP);
kinase A fosforyleert fosforylasekinase, dat glycogeenfosforylase
fosforyleert (uit: fosfatasen); fosforylase maakt glucose-1-fosfaat vrij uit
glycogeen, dat wordt omgezet en als glucose uitgevoerd.
\end{solution}

\begin{solution}{exo:b2:cell-signalling:3}
Oppervlaktereceptoren binden wateroplosbare boodschappers aan het membraan en
werken in seconden via \omterm{prop:b2:cell-signalling:gpcr}{secundaire boodschappers} en kinasen, waarbij ze de
activiteit van bestaande eiwitten veranderen. \omterm{prop:b2:cell-signalling:nuclear}{Kernreceptoren} binden
vetoplosbare boodschappers in de cel en werken in uren als
\omterm{prop:b2:cell-differentiation:transcriptional}{transcriptiefactoren}, waarbij ze veranderen welke eiwitten worden gemaakt.
\end{solution}

\begin{solution}{exo:b2:cell-signalling:4}
\omterm{prop:b2:cell-signalling:gpcr}{Cyclisch AMP}: gemaakt door adenylaatcyclase, afgebroken door
fosfodiësterase. Calcium: komt binnen via kanalen of verlaat het
endoplasmatisch reticulum als antwoord op inositoltrisfosfaat, verwijderd door
pompen. Inositoltrisfosfaat (met diacylglycerol): door fosfolipase C gemaakt
uit een membraanlipide, verwijderd door fosfatasen.
\end{solution}

\begin{solution}{exo:b2:cell-signalling:5}
$\theta = [L]/(K_d + [L])$: $0.09$, $0.5$, $0.91$, $0.99$. Verdubbeling van
de receptoren laat de fractie ongewijzigd en verdubbelt het aantal bezette
receptoren --- en daarmee de respons.
\end{solution}

\begin{solution}{exo:b2:cell-signalling:6}
$50\times 200\times 100\times 300 = 3\times 10^{8}$. Honderd bezette
receptoren geven $3\times 10^{10}$ productmoleculen per seconde --- een
bovengrens, omdat de stadia verzadigen.
\end{solution}

\begin{solution}{exo:b2:cell-signalling:7}
Het vastgezette G-eiwit houdt het cyclase aan, het cAMP blijft hoog, kinase A
houdt het chloridekanaal van de darmcellen gefosforyleerd en open, chloride
wordt uitgescheiden, en natrium en water volgen het het lumen in --- liters
per dag. De vastzetting is een covalente modificatie van het G-eiwit, die pas
ongedaan wordt gemaakt wanneer de cellen zelf worden vervangen, dagen later.
\end{solution}

\begin{solution}{exo:b2:cell-signalling:8}
$0.9\times 10^{-6}\,\text{mol/L}\times 2\times 10^{-12}\,\text{L} =
1.8\times 10^{-18}$ mol, $1.1\times 10^{6}$ ionen (in feite meer, omdat
buffers het meeste van wat binnenkomt binden). Eén kanaal met $10^{6}$ per
seconde zou er een seconde over doen; duizend kanalen doen het in een
milliseconde, de tijdschaal die een signaal nodig heeft.
\end{solution}

\begin{solution}{exo:b2:cell-signalling:9}
De twee routes komen samen op dezelfde enzymen: het kinase A van glucagon
fosforyleert glycogeensynthase (uit) en fosforylasekinase (aan); de cascade
van insuline activeert het fosfatase dat de fosfaatgroepen verwijdert. De
toestand van de enzymen wordt bepaald door het evenwicht tussen kinase- en
fosfataseactiviteit, zodat de verhouding van de twee hormonen telt; als beide
hoog zijn, wint meestal het fosfatase en wordt glycogeen opgeslagen.
\end{solution}

\begin{solution}{exo:b2:cell-signalling:10}
Adrenaline: cAMP diffundeert in minder dan een seconde door een cel van
\qty{20}{\micro m} ($x^{2}/2D$ met $D \approx \qty{3e-10}{m^2/s}$), en elk
enzym van de cascade bestaat al --- seconden. Cortisol: een transcript kost
\qty{60}{s}, een eiwit \qty{100}{s}; uit enkele tientallen mRNA's met elk een
twaalftal ribosomen maakt de cel enkele honderden enzymmoleculen per minuut,
zodat de duizenden die voor een meetbaar effect nodig zijn uren vergen, na
de minuten die cortisol nodig heeft om binnen te komen en het DNA te bereiken.
\end{solution}

\begin{solution}{exo:b2:cell-signalling:11}
De receptor geeft voortdurend een signaal: de MAP-kinasecascade drijft de cel
tot deling zonder enige groeifactor, en de cel negeert de afwezigheid van het
signaal dat haar normaal afremt. Omdat één \omterm{def:b2:genome-diversification:mutation}{mutatie} een controle op de deling
wegneemt, behoren zulke receptoren (en de kinasen stroomafwaarts) tot de
meest voorkomende oncogenen. Een middel dat de ATP-plaats van het kinase
bezet, stopt de fosforyleringen en legt de receptor stil --- het principe van
verscheidene kankermiddelen.
\end{solution}

\begin{solution}{exo:b2:cell-signalling:12}
Endocrien: bereikt elke cel in een minuut, goedkoop per cel, kan geen cel
afzonderlijk aanspreken, werkt minuten tot dagen --- onmisbaar voor toestanden
van het hele lichaam (vasten, groei, voortplanting). Zenuw: milliseconde, één
doelwit, dure bedrading, werkt milliseconden --- onmisbaar voor beweging en
snelle reflexen. Beide worden voor dezelfde boodschap gebruikt waar zowel
snelheid als bereik gewenst zijn: de sympathische zenuwen en het bijniermerg
leveren bij een schrik allebei noradrenaline of adrenaline.
\end{solution}

\begin{solution}{pb:b2:cell-signalling:1}
\textbf{1.} $1\,\text{nmol}/5\,\text{L} = \qty{0.2}{nmol/L}$.
\textbf{2.} $\theta = 0.2/1.2 = 0.17$: ongeveer 3300 bezette receptoren.
\textbf{3.} Drie halveringstijden: \qty{0.025}{nmol/L}; $\theta = 0.024$.
\textbf{4.} Ongeveer een halve tot een hele minuut, de tijd die het \omterm{def:b2:blood-circulation:blood}{bloed}
nodig heeft om rond te gaan; de zenuw levert haar transmitter binnen een
seconde rechtstreeks aan de lever.
\textbf{5.} Dezelfde fractie, maar $33\,000$ bezette receptoren: een tien keer
zo grote respons.
\textbf{6.} Bij \qty{1}{nmol/L} is een receptor met $K_d = \qty{1}{\micro
mol/L}$ voor een duizendste bezet: geen respons. Door $K_d$ af te stemmen op
het circulerende bereik ligt het steile deel van de bezettingscurve waar het
\omterm{def:b2:cell-signalling:kinds}{hormoon} werkelijk varieert.
\textbf{7.} $20\times 100\times 1\times 50\times 50\times 1000 =
5\times 10^{9}$ glucosemoleculen per seconde per bezette receptor.
\textbf{8.} $3300\times 20\times 100 = 6.6\times 10^{6}$ cAMP in de eerste
seconde; in \qty{5e-12}{L}, $1.1\times 10^{-17}$ mol, ongeveer
\qty{2}{\micro mol/L}.
\textbf{9.} $3300\times 5\times 10^{9} = 1.7\times 10^{13}$ per seconde, $2\times
10^{13}$ per minuut per cel.
\textbf{10.} $2\times 10^{11}$ cellen $\times 10^{15} = 2\times
10^{26}$ moleculen per minuut, 330 mol, \qty{6e4}{g} per minuut ---
tienduizend keer de waargenomen \qty{5}{g/min}. De versterkingen
vermenigvuldigen zich alleen zolang geen enkel stadium verzadigd is; in
werkelijkheid wordt het fosforylase beperkt door zijn substraat en door de
fosfatasen, en de uitvoer van glucose door transporters. De versterking van
de cascade wordt besteed aan snelheid en gevoeligheid, niet aan productie.
\textbf{11.} Bij \qty{5}{g/min} twintig minuten; de schrik is in minuten
voorbij en de uitschakelaars stoppen de cascade lang daarvoor.
\textbf{12.} Hydrolyse van GTP door het G-eiwit (seconden);
fosfodiësterase op cAMP (seconden); fosfatasen op de gefosforyleerde enzymen
(seconden tot minuten); \omterm{prop:b2:reproductive-hormones:pulses}{desensitisatie van de receptor} (minuten). Het
snelst zijn de eigen klok van het G-eiwit en het fosfodiësterase.
\textbf{13.} cAMP wordt niet afgebroken: het hoopt zich op en blijft, het
kinase blijft aan, en de glucoseproductie is groter en duurt langer --- de
alertheid van koffie.
\textbf{14.} $4000/50 = \qty{80}{s}$.
\textbf{15.} $600/5 = \qty{120}{s}$ per enzym per ribosoom; 30 per
ribosoom per uur, 600 per mRNA per uur.
\textbf{16.} $100\times 600 = 60\,000$ enzymen per uur; $10^{5}$ na
ongeveer \qty{1.7}{h}.
\textbf{17.} Tien tot twintig minuten voor cortisol om binnen te komen, te
binden en zijn genen te bereiken, een uur voordat de eerste mRNA's talrijk
zijn en worden getranslateerd, en daarna twee uur opstapeling: ongeveer vier
uur.
\textbf{18.} Een cascade verandert de activiteit van enzymen die bestaan;
de taak van cortisol is dagenlang te veranderen welke enzymen bestaan, en dat
kan alleen transcriptie. Een schrik kan niet vier uur op nieuwe enzymen
wachten.
\textbf{19.} Meer receptoren betekent meer bezette receptoren bij elke
adrenalineconcentratie: de bezettingsfractie is onveranderd, maar de respons
bij elke concentratie is groter --- de curve wordt opgeschaald, niet
verschoven. Een traag \omterm{def:b2:cell-signalling:kinds}{hormoon} stelt de versterking van een snel \omterm{def:b2:cell-signalling:kinds}{hormoon} in.
\textbf{20.} Glucose komt de $\beta$-cel binnen en wordt gemetaboliseerd;
het ATP stijgt en sluit een ATP-gevoelig kaliumkanaal; het membraan
depolariseert; spanningsafhankelijke calciumkanalen gaan open; calcium zet de
exocytose van insulinekorrels in gang --- metabolisme dat via een elektrische
stap aan secretie is gekoppeld.
\textbf{21.} De twee receptoren starten twee routes die samenkomen op dezelfde
enzymen: het kinase A van glucagon fosforyleert ze, de cascade van insuline
activeert het fosfatase dat ze defosforyleert; de toestand van de enzymen, en
dus de richting van het glycogeenmetabolisme, volgt de verhouding van de
twee.
\textbf{22.} Nuchtere glucose hoog (de lever blijft haar ongeremd maken),
glucose na een maaltijd zeer hoog (geen opname door spier- en vetweefsel);
zonder insuline geven de vetvoorraden vetzuren af, die de lever omzet in
ketonzuren --- ketoacidose.
\textbf{23.} Werkende spieren nemen glucose op zonder insuline, zodat de
glucose tijdens een lange duurloop daalt terwijl de ingespoten insuline,
ongeregeld, de lever blijft onderdrukken: hypoglykemie en ineenstorting.
\textbf{24.} Glucoselus: als sensor de $\beta$- en $\alpha$-cellen, twee
effectorhormonen voor de twee richtingen, een hoge versterking die het niveau
tot op een vijfde vasthoudt, minuten om te werken. \omterm{def:b2:blood-pressure:baroreflex}{Baroreflex}: rekreceptoren,
het \omterm{def:b2:heart:anatomy}{hart} en de vaten als effectoren, versterking ongeveer vier, een seconde
om te werken.
\textbf{25.} Bezetting $0.17$ na de afgifte; versterking $5\times 10^{9}$ per seconde
per bezette receptor; productie van de lever in werkelijkheid
\qty{5}{g/min}; cortisol heeft ongeveer vier uur nodig.
\end{solution}
